\section{The Laurent model for analytic compact supports}
\label{sec:laurent}

We now express the compact supports of the geometric calculation as
convergent Laurent series. The output is an explicit module with the
actual stabilizer action and finite cohomology and homology. Its
finite-field form will carry the residue map that is transpose to the
Cartier transfer, including series with unbounded root denominators.

Fix $1\leq m<n$, put $r=n-m$ and $s=r-1$, and retain the
whole finite free model
\[
 A=\kk[[x,y_1,\ldots,y_s]],\qquad B=A[1/x],\qquad
 A'=\bigoplus_{\nu=1}^{N}A b_\nu,\qquad B_U=eA'[1/x]
\]
of \cref{prop:mark-whole-model}.  Thus $A'[1/x]$ is finite
\'{e}tale over $B$, its multiplication laws in the displayed basis are
integral, and $e$ is a $G$-invariant idempotent.  All operations involving
$e$ below take place in the whole free module.  In particular no
specialization of $e$ at $x=0$ is used.

Write $F=\Cflat$, fix its absolute value, and let $X$ be the perfected
characteristic-$p$ analytic domain
\[
 0<|x|<1,\qquad |y_j|<1\quad(1\leq j\leq s)
\]
over $F$.  Its finite marked cover is denoted by $Z_U$.  The sheaf
$\cO^\flat$ on this characteristic-$p$ space is its structural sheaf.
The present calculation consequently concerns a different space from the
untilted affine line in \cref{prop:kum-line}.

\subsection{Complete Laurent steps}

For a compact profinite set $\mathcal D$, put
$R_{\mathcal D}=C(\mathcal D,F)$, with its supremum norm.  An empty
parameter set will not be used; when there is no parameter we take a
one-point set.  Set
\[
 \mathcal E=\Z[1/p]\times\bigl(\Z[1/p]_{>0}\bigr)^s,
 \qquad |J|=J_1+\cdots+J_s.
\]
If $s=0$, the second factor has one element and $|J|=0$.
For $L>\delta>0$ define
\begin{equation}\label{eq:laurent-weight}
 w_{L,\delta}(I,J)=
 \begin{cases}
 LI-\delta|J|,&I\geq0,\\
 \delta I-\delta|J|,&I<0.
 \end{cases}
\end{equation}
Let $\mathscr L_{L,\delta}(\mathcal D)$ consist of the sums
\[
 f=\sum_{(I,J)\in\mathcal E}a_{I,J}x^Iy^{-J},
 \qquad a_{I,J}\in R_{\mathcal D},
\]
such that the numbers
$\|a_{I,J}\|_{\sup}\exp(-w_{L,\delta}(I,J))$ tend to zero outside
finite subsets of $\mathcal E$.  Equip this space with the maximum of
these numbers as norm.  It is a complete nonarchimedean normed module:
coefficientwise limits of a Cauchy sequence have uniformly small tails,
which proves both the required decay and convergence in the norm.
Finite sums are dense.  Define
\[
 \mathscr L(\mathcal D)=\colim_{L,\delta}
       \mathscr L_{L,\delta}(\mathcal D).
\]
The transitions increase $L$ and decrease $\delta$.  They are norm
nonincreasing, since these changes increase every weight.  Rational pairs
$L>\delta>0$, or the sequence $(L,\delta)=(j,j^{-1})$ with $j\geq2$,
are cofinal.  The colimit here and below has the complete-step convention
of \cref{sec:continuous}; a parameter family belongs to one step.

Multiplication by an element of $B$ is multiplication of Laurent series
followed by discarding every monomial whose exponent in some $y_j$ is
nonnegative.  This prescription is well defined in the complete steps.
Indeed, if $a\in x^{-c}A$, every monomial of $a$ has exponents
$(A_0,B_0)$ with $A_0\geq-c$ and all $(B_0)_j\geq0$.  Whenever its
product survives the projection,
\begin{equation}\label{eq:laurent-matrix-bound}
 w_{L,\delta}(I+A_0,J-B_0)
       \geq w_{L,\delta}(I,J)-Lc.
\end{equation}
This follows because the increasing piecewise linear function
$I\mapsto LI$ for $I\geq0$, $I\mapsto\delta I$ for $I<0$, has slopes
between $\delta$ and $L$.  At a fixed root grid, multiplying one
monomial by $a$ gives a convergent series: after projection only finitely
many $y$-exponents survive below the original pole orders, and the
remaining $x$-orders tend to infinity.  The bound extends multiplication
from finite sums to arbitrary sums, proves its continuity, and justifies
associativity by dense finite-sum approximation.  A matrix with entries
in $x^{-c}A$ has the same bound, up to the maximum norm of its finitely
many entries.

Consequently $e$ is a continuous chain-compatible idempotent on the
whole ambient Laurent module.  We use the notation
\[
 \mathscr W_{\mathcal D}
   =e\bigl(\mathscr L(\mathcal D)\otimes_A A'\bigr),
 \qquad
 W_{\mathcal D}
   =e\bigl(\mathscr L(\mathcal D;\kk)\otimes_A A'\bigr),
 \qquad W=W_{\{*\}},
\]
where $\mathscr L(\mathcal D;\kk)$ denotes the same sums with
coefficients in $C(\mathcal D,\kk)$.  The tensor notation means the
finite sum on the displayed basis, with the multiplication action just
defined; it introduces no additional completion.  The image of $e$ in
each step is closed, since it is the kernel of $1-e$, and is therefore
complete.

For finite-field coefficients the decay condition has a useful strong
consequence.  For every real $M$, a family in one step has only finitely
many nonzero coefficients with $I\leq M$, uniformly on $\mathcal D$.
Indeed, every nonzero coefficient has supremum norm one, whereas
\[
 I\leq M\quad\Longrightarrow\quad
 w_{L,\delta}(I,J)\leq L\max(M,0).
\]
Infinitely many such coefficients would contradict decay.  Each of the
remaining finitely many coefficient functions is locally constant.
In particular bounded-$x$ truncation produces a locally constant family
of finite principal parts, even when the original exponents have
unbounded $p$-power denominators.  For a one-point parameter the steps
are separable: their finite sums have coefficients in the finite field
and indices in the countable set $\mathcal E$.  Thus the complete steps
of $W$ are Polish additive groups.  A condition that coefficients vanish
outside a prescribed subset, for example $I\geq a>0$, is closed on each
such step.

\subsection{The support complex and coefficient descent}

Let $Q_k=|\kk|$.  The operator $\varphi$ used in this subsection raises
coefficients in $F$ to their $Q_k$th powers and fixes $x^I$, $y^{-J}$,
and every $b_\nu$.  It fixes the structure constants of $A'$ and the
matrix of $e$, since they are defined over $\kk$.  This is
\emph{coefficient-only} Frobenius.  For a complex $K$ with this operator
write
\[
 K^{\varphi=1}=\operatorname{fib}(\varphi-1:K\longrightarrow K).
\]

\begin{proposition}[Laurent support comparison]\label[proposition]{prop:laurent-model}
With the uniform compact-support convention of
\cref{lem:tower-supports}, there are natural quasi-isomorphisms
\begin{equation}\label{eq:laurent-descent}
 \RGamma_c(Z_U,\cO^\flat;\mathcal D)
       \simeq\mathscr W_{\mathcal D}[-r],
 \qquad
 \RGamma_c(Z_U,\cO^\flat;\mathcal D)^{\varphi=1}
       \simeq W_{\mathcal D}[-r].
\end{equation}
They are compatible with finite marking maps, the idempotent $e$, and
the actions defined over $\kk$.  The support cohomology is therefore
concentrated in cohomological degree $r$, both before and after
coefficient descent.  Every fixed normalized root basis
$b_1^{1/h},\ldots,b_N^{1/h}$ gives the same module and the same
complete-step category.
\end{proposition}

\begin{proof}
We first calculate on $X$ with scalar coefficients.  Compact bounds
are cofinal among boxes
\[
 K(a,b,\rho)=
 \{a\leq|x|\leq b,\ |y_j|\leq\rho\ (1\leq j\leq s)\},
 \qquad 0<a<b<\rho<1.
\]
In fact a quasicompact bound has its $x$-norm bounded away from zero and
all coordinate norms bounded strictly below one; one can then choose
$\rho$ larger than its $x$-bound and every $y$-bound.  Interleaving a
slightly smaller closed box and a slightly larger open box gives the
same support colimit.  This permits strict gaps in every radius used
below and avoids assertions about a particular higher-rank boundary
point.

The complement of a box is the union of the two $x$-exteriors and the
$s$ individual $y$-exteriors.  Its finite \v Cech restriction diagram,
augmented by sections on $X$, calculates the support fiber.
Every intersection is a product of open discs, punctured discs, and
annuli.  Exhaust it by a countable sequence of closed polyannuli and
polydiscs, after adjoining all coordinate roots. The root polydisc,
rational-localization and product constructions make these closed domains
affinoid perfectoid
\citep[Proposition~5.20, Theorem~6.3\textup{(i)--(ii)}, and Proposition~6.18]{scholzePerfectoid2012}.
Their structural sheaf is v-acyclic by
\cite[Proposition~8.8]{scholzeDiamonds2017}.  Their function algebras
are the corresponding completed Laurent algebras.  The restriction
maps have dense image: finite Laurent polynomials belong to both
domains and are dense in the smaller-domain algebra.  With parameters,
finite sums whose coefficients are in $R_{\mathcal D}$ have the same
property.  Equivalently, first approximate a continuous parameter
family by a finite clopen partition and then approximate its finitely
many values by Laurent polynomials.  The complete normed-space
inverse-limit lemma, \cref{lem:kum-banach-limits}, eliminates the first
derived inverse limit.  There are no higher derived limits for a
countable sequence.  Thus these open intersections have their ordinary
Laurent functions in degree zero, with no additional cohomology.

Here is the resulting splitting of the support diagram.  In an
unpunctured coordinate $y_j$, a function on the exterior annulus is the
sum of its nonnegative and strictly negative Laurent parts.  The
nonnegative part extends to the whole disc, since it converges at every
radius less than one.  Its summand in the restriction complex is
contractible.  The negative part is the quotient summand, in degree
one.  These projections have norm at most one on each closed annulus;
on open domains one first inserts a radius gap.  They are
$R_{\mathcal D}$-linear and act coefficientwise in the other
coordinates.

For the punctured coordinate write an exterior pair as
$(f_{\rm in},f_{\rm out})$.  The negative part of $f_{\rm in}$ extends
to the whole punctured disc: convergence towards zero is the only
extra condition for a negative-power series.  Likewise the
nonnegative part of $f_{\rm out}$ extends to the whole punctured disc.
Subtracting these global functions diagonally leaves the unique pair
\[
 \bigl(f_{\rm in,+}-f_{\rm out,+},\,
       f_{\rm out,-}-f_{\rm in,-}\bigr).
\]
Encode it by the single Laurent series
\[
  f=(f_{\rm in,+}-f_{\rm out,+})
       -(f_{\rm out,-}-f_{\rm in,-}).
\]
Thus its normal representative is $(f_+,-f_-)$.  The minus sign
ensures that multiplication or substitution crossing the exponent
$I=0$ is represented by the usual Laurent operation: eliminating an
outer nonnegative part by a diagonal section introduces its negative
at the inner end.  This quotient is again in degree one, and it
retains every exponent $I\in\Z[1/p]$, including zero exactly once.

Apply these splittings successively to the finite product restriction
diagram.  A summand which extends in any one coordinate is
contractible by the corresponding one-coordinate splitting.  The only
remaining summand has every $y$-exponent negative and the indicated
two-ended $x$-representative.  It lies in degree $s+1=r$.
The splittings commute in distinct coordinates, so this is a
calculation of the entire complex, not only of its highest quotient.
It also proves the assertion for $s=0$.

For the surviving inner positive $x$-powers choose a radius
$\exp(-L)$ slightly below the inner support radius.  For the surviving
negative $x$-powers and negative $y$-powers choose a common outer
radius $\exp(-\delta)$ slightly above every outer support radius.
Their Gauss norms are exactly
\[
 \|a_{I,J}\|_{\sup}\exp(-LI+\delta|J|)
 \quad(I\geq0),\qquad
 \|a_{I,J}\|_{\sup}\exp(-\delta I+\delta|J|)
 \quad(I<0).
\]
Convergence on a closed polyannulus is precisely decay outside finite
subsets for these norms.  Every support class has these bounds after
inserting the radius gaps.  Conversely a series in one such step
converges on the smaller inner domain and the larger outer collars,
and hence supplies a class for a box with separated radii.  Enlarging
supports decreases the inner radius and increases the outer radii;
on these representatives the map retains the same coefficients.
The support colimit is therefore exactly $\mathscr L(\mathcal D)$.
This calculation uses the finite restriction diagram followed by a
filtered colimit of supports.  It does not interchange that colimit
with an inverse limit of shrinking supports.

We next pass to the actual finite marked algebra.  On a closed base
chart with perfectoid algebra $R$, the whole algebra
\[
 A'[1/x]\otimes_B R=\bigoplus_{\nu=1}^N Rb_\nu
\]
is already complete as a finite free $R$-module.  It is finite
\'{e}tale and affinoid perfectoid by
\cite[Theorem~6.1(i)]{scholzeDiamonds2017}.  Relative Frobenius for an
\'{e}tale algebra is an isomorphism
\cite[Lemma~41.14.3, Tag~0EBS]{stacks}; consequently this algebra is
also the perfected marked extension, rather than only an
unperfected approximation.  Restrictions use the same whole basis,
so the function complex is the finite sum of the scalar complexes.
The matrix $e$ commutes with every restriction and is an actual chain
idempotent.  Taking its image commutes with the finite fibers and
filtered support colimit.  For the intermediate inverse limits its
image remains complete and the transition maps remain dense: apply
the target idempotent to ambient approximants.  Hence the preceding
inverse-limit argument applies there as well.

This computation uses all compact supports on $Z_U$.  Pullbacks of
quasicompact bounds under a finite map are quasicompact, and the image
of a quasicompact bound is a quasicompact bounded subset of $X$.
Thus pulled-back boxes form a cofinal support system upstairs.  This
establishes the first quasi-isomorphism in
\cref{eq:laurent-descent} on the actual cover.

Coefficient Frobenius and its inverse are continuous between steps.
For example, raising a coefficient to its $Q_k$th power changes a
weight $w$ to $Q_kw$; this is the step with parameters
$(Q_kL,Q_k\delta)$, followed if needed by a cofinal enlargement.
It remains to check surjectivity of $\varphi-1$, including parameters.
For $b\in F$, the equation
\[
 a^{Q_k}-a=b
\]
has a solution with $|a|\leq|b|$.  If $|b|<1$, take
$a=-\sum_{j\geq0}b^{Q_k^j}$.  If $|b|=1$, any root is integral.
If $|b|>1$, a root has norm $|b|^{1/Q_k}$.  These choices have
continuous local branches: around a root $a_0$ for $b_0$, add the
small solution for $b-b_0$.  Near a point with $|b_0|\geq1$ take
$|b-b_0|<1$, so the same bound is retained; on $|b|<1$ use the
displayed series itself.  For a continuous function on a compact
profinite $\mathcal D$, a finite clopen refinement of these local
charts therefore gives a continuous solution with
$\|a\|_{\sup}\leq\|b\|_{\sup}$.

Apply this construction separately to every Laurent coefficient.
The weighted norms do not increase, so all the chosen coefficients
assemble to a family in the original complete step.  Uniformly small
tails make that assembled family continuous even though the finite
clopen partitions used for different coefficients may differ.
Solving on the whole finite sum and then applying $e$ proves
surjectivity on $\mathscr W_{\mathcal D}$.  Its kernel is exactly
$W_{\mathcal D}$, since the fixed elements of $F$ are $\kk$ and $e$
is fixed.  This proves the second quasi-isomorphism.

Finally, relative Frobenius identifies each fixed root version of
$A'[1/x]$ with its base change to $B^{1/h}$.  Thus
$b_\nu^{1/h}$ is a basis over the perfected base.  Any two fixed
such bases are related by invertible matrices over a finite root
version of $B$, with bounded $x$-poles and no negative $y$-powers.
The bound \cref{eq:laurent-matrix-bound}, rescaled to that root grid,
makes both basis changes continuous.  If a fixed matrix before
taking roots has pole bound $c$, its $h$-root has bound $c/h$.
Moreover, when $h\mid h'$, each $b_\nu^{1/h}$ expands integrally
in the basis $b_\mu^{1/h'}$: it is the $(h'/h)$th power of a basis
element, and the multiplication laws of the root model are integral.
These observations prove both the basis assertion and compatibility
with the stated operations.
\end{proof}

\subsection{Uniform continuity of the group action}

\begin{lemma}\label[lemma]{lem:laurent-action}
The $G$-action on $W$ and on $\mathscr W_{\mathcal D}$ is jointly
continuous from a complete step, with a compact set of group and
parameter values, into a complete step.  It is uniformly
equicontinuous there.  The same holds for multiplication by the
finite marked function algebra, fixed root-basis changes, and the
finite matrices occurring in continuous group resolutions.
Absolute Frobenius is an invertible additive operator on these
modules.  Its sufficiently divisible powers fixing $\kk$ commute
with $G$ and act by chain automorphisms on the group complexes over
$\kk$.
\end{lemma}

\begin{proof}
We check the scalar coordinate changes first.  An integral formal
automorphism preserving $(x)$ factors into three types:
\[
 x\longmapsto xu(x,y),\quad y\longmapsto y;
 \qquad x\longmapsto x,\quad y\longmapsto y+c(x);
 \qquad x\longmapsto x,\quad y\longmapsto v(x,y),
\]
where $u$ is an integral unit, $c(x)\in x\kk[[x]]^s$, and the last
map preserves $(y)$ and has invertible linear Jacobian in $y$ over
$\kk[[x]]$.  To obtain the factorization, first remove the $x$-unit
substitution by its formal inverse.  For the residual map fixing
$x$, subtract the value of its $y$-coordinates at $y=0$; the
remaining map fixes $(y)$ and its $y$-Jacobian is invertible.
All inverse maps are integral.  These operations are continuous on
formal jets, uniformly in a compact family.

Fix a root grid $1/h$ and a monomial $x^Iy^{-J}$.  For each of the
three types every surviving output monomial $x^{I'}y^{-J'}$ satisfies
\begin{equation}\label{eq:laurent-action-order}
 I'\geq I,\qquad |J'|\leq |J|+(I'-I).
\end{equation}
For the $x$-unit substitution, the expansion of $u^I$ on the root
grid has only nonnegative additional $x$- and $y$-orders, giving
the stronger bound $|J'|\leq|J|$.  For translation, write
$X=x^{1/h}$ and $Y_j=y_j^{1/h}$.  Expanding
$(Y_j+c_j(x)^{1/h})^{-hJ_j}$, an increase $a_j/h$ in its pole
requires the factor $c_j(x)^{a_j/h}$, whose $x$-order is at least
$a_j/h$.  Summing over $j$ gives
\cref{eq:laurent-action-order}.

For the third type use finite principal-parts duality, which also
avoids choosing a preferred expansion of a general linear
combination of the $y_j$.  On the grid just fixed the negative
$Y$-powers form
\[
 H^s_{(Y)}\bigl(\kk[[X,Y_1,\ldots,Y_s]]\bigr),
\]
with coefficients extended in $X$ as needed.  A principal part
$Y^{-j}$ pairs with $Y^{j-1}\,dY_1\cdots dY_s$ by coefficient
extraction.  This is a perfect pairing between every finite
principal-part truncation and the corresponding quotient of the
power-series ring.  If $v$ preserves $(Y)$, then its inverse and
the determinant of its inverse Jacobian are integral, and
\[
 (v^{-1})^*(Y^{j'-1}\,dY_1\cdots dY_s)
\]
has total $Y$-degree at least $|j'|-s$.  It cannot pair nontrivially
with $Y^{-j}$ when $|j'|>|j|$.  Every coefficient remains in
$\kk[[X]]$.  Thus here $I'\geq I$ and $|J'|\leq|J|$.
The pairing description is the usual change-of-variables map on
the finite \v Cech complex: it is first checked on monomial
fractions by coefficient extraction and then on their finite
principal-part span.  Hence it represents the geometric support
pullback.  The case $s=0$ requires no such step.

The two inequalities in \cref{eq:laurent-action-order} are stable
under composition.  They imply
\[
 w_{L,\delta}(I',J')\geq w_{L,\delta}(I,J)
 \qquad(L>\delta>0).
\]
When both $x$-exponents are negative this uses the slope $\delta$;
when both are nonnegative it uses $L\geq\delta$; when they cross
zero the excess is at least $(L-\delta)I'$.  The output of a fixed
monomial is convergent.  For translations, after $I'$ becomes
nonnegative the weight increases by at least $L-\delta$ per
additional unit of pole, and at a fixed grid only finitely many
terms precede that crossing.  For a map preserving $(y)$, the
total $y$-pole order is bounded and its remaining $x$-series
converges.  The $x$-unit case has the same property.

These expansions also occur on actual analytic collars.  For
example
\[
 (y+x)^{-1}=\sum_{j\geq0}(-x)^jy^{-j-1}
\]
converges on $|x|<|y|$, and its term weights are
$(L-\delta)j-\delta$.  Any compact support bound can be enlarged
so that its $y$-radius exceeds its $x$-radius.  The same enlargement
works for every integral $c(x)$ because $|c(x)|\leq|x|$.
For maps preserving $(y)$ the ideal and its inverse have the same
total-radius bounds; finite \v Cech refinements give the
principal-parts pullback above.  Thus the algebraic formulas agree
with the canonical support action used in
\cref{prop:laurent-model}.

All coefficients of these integral maps have norm at most one.
The weight estimate therefore gives a uniform operator bound on
each complete scalar step.  For a fixed monomial, every
bounded-weight output truncation involves finitely many monomials
on its fixed root grid and depends on finitely many formal jets.
It consequently varies continuously with the group element.
For an arbitrary convergent sum, discard a uniformly small tail
and apply this assertion to its finite remainder.  This proves
joint continuity and equicontinuity, including all compact
parameters.  Operator-norm continuity in the group element is
neither asserted nor needed.

On the whole finite model, the group action is the scalar pullback
just treated followed by a finite basis matrix.  The integral
model of \cref{prop:mark-whole-model} and compactness give a uniform
pole bound for those matrices and their inverses.  Their entries
vary continuously in the bounded-lattice topology.  Apply
\cref{eq:laurent-matrix-bound}; then pass to the invariant image
of $e$.  This proves the same assertions for the marked module.
That bound also proves the claims about multiplication and root
bases.

For completeness, these estimates allow integration against
compact free resolution terms as in \cref{lem:cts-comparison}.
If a compact family has image $K_0$ in a complete linearly
topologized $\kk$-step, the closed $\kk$-linear span of $K_0$ is
compact.  Modulo any open linear subgroup its image is the span
of a finite set over the finite field, hence finite; this proves
total boundedness, and completeness gives compactness.  Enlarge
the step first to contain the compact group translates if
necessary.  Summation against compact distributions, and the
resolution comparison homotopies, therefore remain inside a
complete step.  No boundary image is replaced by its closure.

Finally, absolute Frobenius raises coefficients and basis
elements to their $p$th powers \emph{and scales every exponent by
$p$}.  The relative-Frobenius basis identifications in
\cref{prop:laurent-model} and the rescaled weight condition show
that this operation and its inverse act continuously between
steps.  It commutes with the geometric action.  A power fixing
$\kk$ is $\kk$-linear, so it commutes with the differentials and
homotopies of resolutions over $\kk$.  This operation is distinct
from the coefficient-only $\varphi$ in
\cref{eq:laurent-descent}.
\end{proof}

\subsection{Finiteness of cohomology and homology}

\begin{corollary}\label[corollary]{cor:laurent-finite}
At every sufficiently deep product marking level used in
\cref{cor:ind-finite}, both $H^a_{\cts}(G,W)$ and
$H_i^{\cts}(G,W)$ are finite-dimensional over $\kk$ for every
$a,i\geq0$.  Since $\kk$ is finite, these are finite groups.
\end{corollary}

\begin{proof}
Put $K=\RGamma(G,\RGamma_c(Z_U,\cO^\flat))$, with no external
parameter.  By \cref{cor:ind-finite}, every $H^j(K)$ is
finite-dimensional over $F$.  By \cref{lem:laurent-action} the
complete-step comparisons in \cref{lem:cts-comparison} apply.
The finite-type resolution in each degree commutes with taking
the fiber of $\varphi-1$.  Thus
\begin{equation}\label{eq:laurent-group-descent}
 K^{\varphi=1}\simeq\RGamma(G,W)[-r].
\end{equation}
One can equally form this fiber on bounded-step bar cochains:
the parameter-uniform surjectivity proved above applies to the
compact bar parameters, and the two descriptions agree.

We recall explicitly the finite-dimensional semilinear argument.
If $V$ is an $F$-vector space of dimension $d$ and $\psi$ is an
invertible $Q_k$-semilinear map, choose coordinates
$\psi(z)=Az^{[Q_k]}$ with $A$ invertible.  For any target $t$,
the equations $Az^{[Q_k]}-z=t$ are, after multiplication by
$A^{-1}$, monic equations with pairwise coprime leading
monomials $z_1^{Q_k},\ldots,z_d^{Q_k}$.  Monomial reduction gives
a quotient algebra of dimension $Q_k^d$, with basis
$\prod z_i^{a_i}$ for $0\leq a_i<Q_k$.  Their Jacobian is the
invertible matrix $-A^{-1}$, so that algebra is finite \'{e}tale.
As $F$ is algebraically closed, every fiber has $Q_k^d$ points.
Consequently $\psi-1$ is surjective and its kernel is a
$d$-dimensional $\kk$-vector space.

Apply this to the bijective semilinear action of $\varphi$ on
each $H^j(K)$.  The long exact sequence of its fiber and the
surjectivity in the preceding degree give
\[
 H^j(K^{\varphi=1})=
       H^j(K)^{\varphi=1},
\]
which is finite over $\kk$.  The shift in
\cref{eq:laurent-group-descent} gives finiteness of
$H^a_{\cts}(G,W)$.

The marked algebra acts continuously on $W$ by
\cref{eq:laurent-matrix-bound}, and its finite subgroups have
the trace-one functions supplied by
\cref{prop:mark-finite-isotropy}.  All the hypotheses of
\cref{prop:cts-full-duality} are therefore satisfied.  Its
orientation for $G=P_n$ is trivial, and it identifies
\[
 H_i^{\cts}(G,W)\simeq
 H^{\dim(G)-i}_{\cts}(G,W).
\]
The right side is finite (and zero in negative degree), proving
the assertion.  The argument uses parameter families to justify
the complexes; it makes no claim of finite dimension over $F$
or $\kk$ with an infinite external parameter retained.  In that
case the respective scalar rings are $R_{\mathcal D}$ and
$C(\mathcal D,\kk)$.
\end{proof}
